\section{Implementation of the method in numerical simulations}
\label{sec:duef}

To implement the ideas of the previous section,
let us consider the generic local two-fermion operator
\beq O_{\Gamma}(x) = \bar{\psi}(x)\Gamma \psi(x) \label{eq:2ferm} \eeq
where $\Gamma$ is any Dirac matrix.
The lattice, bare Green function of $O_{\Gamma}$, between off-shell quark
states,
can be obtained from
\beq
G_{O}( x, y )=\langle \psi(x) O_{\Gamma}(0)\bar{\psi}(y) \rangle=
\frac{1}{N}\sum^{N}_{i=1} S_i(x \vert 0) \Gamma S_i(0 \vert y),
\label{eq:gxy} \eeq
where $i=1,N$ labels the gauge field configurations and
$S_i(x \vert 0) $ is the quark propagator on the single configuration $i$,
obtained by inverting the discretized Dirac operator. We denote by
 $S_i(x \vert y)$ the inverse Dirac operator, which is not translationally
invariant, in contrast to $S(x-y)$, which is the usual quark propagator.
    The gauge field configurations
are generated with some standard, gauge-invariant algorithm;
the gauge fields and the quark
propagators are then rotated into the Landau gauge, which
will be  defined in sec.\ \ref{sec:nonper}.
 From eq.\ (\ref{eq:gxy}), using $S_i(x \vert y)= \gamma_5
S_i(y \vert x )^\dagger \gamma_5$, one gets
\beq
G_{O}(pa)&\equiv &\int d^4 x d^4 y  e^{-i \,( p \cdot x- p \cdot y)}
 G_O( x ,  y ) \nn \\
&=&\frac{1}{N} \sum_{i=1}^{N} \Bigl(\int d^4 x S_i(x \vert 0)
e^{-i \, p \cdot x} \Bigr) \Gamma \Bigl(\int d^4 y S_i(0 \vert y)
e^{i \, p \cdot y} \Bigr)
\nn \\ &=&\frac{1}{N} \sum_{i=1}^{N}  S_i(p \vert 0)
\Gamma  \gamma_5 \Bigl(\int d^4 y S^\dagger_i(y \vert 0)
e^{i \, p \cdot y} \Bigr) \gamma_5
\nn \\& =&\frac{1}{N} \sum_{i=1}^{N} S_i(p \vert 0)
 \Gamma  \Bigl( \gamma_5 S^\dagger_i(p \vert  0) \gamma_5
 \Bigr) ,
\label{eq:nuovo}
\eeq
where we have defined
\beq
S_i(p \vert y)=\int d^{4}x S_i(x \vert y)
e^{-i \,( p \cdot x)},
\eeq
because, on a single configuration, the quark propagator
is not translationally invariant. Traslational invariance is recovered
after averaging over the configurations. We can then write
\beq
S(pa)= \langle S(p \vert 0) \rangle= \frac{1}{N}\sum^{N}_{i=1} S_i(p \vert 0)
. \label{eq:nuovo1} \eeq

Eqs.\ (\ref{eq:nuovo}) and (\ref{eq:nuovo1}) give $G_O(pa)$ and
$S(pa)$ in terms of quark propagators, computed by inverting  the Dirac
operator in the  presence of a background gauge field configuration, generated
by numerical simulation. With these quantities at hand, one can
apply the strategy and  the  formulae,
introduced in sec.\ \ref{sec:idea}. We add some information, that can help
the reader in practical applications.
The quark wave-function renormalization constant $Z_\psi$ has been
 computed using
\beq
Z_{\psi}= \Gamma_{V^L} \times Z_{V^L}= \frac{1}{48} {\rm tr}
\left( \Lambda_{V^L_{\mu}} (p a)\gamma_\mu\right)\vert_{p^2=\mu^2}\times
  Z_{V^{L}}
\label{eq:zpsil}
\eeq
where $\Lambda_{V_\mu^L}(pa) $ is defined as in eq.\ (\ref{eq:lam});
the trace
is taken over colours and spins and $Z_{V^L}$ is the renormalization
constant of the local vector current $V^L_\mu(x)=\bar \psi(x)
\gamma_\mu \psi(x)$. Equations (\ref{eq:zpsi})
and (\ref{eq:zpsil}) are two equivalent definitions of
$Z_\psi$. We have used eq.\ (\ref{eq:zpsil}) instead
of eq.\ (\ref{eq:zpsi}), because the former is more convenient when,
as in our case, one has only computed  quark propagators which
stem from the  origin. The  three-point ($q-V-q$) Green function is obtained by
inserting
the local  vector   current in the origin.    $Z_{V^L}$ can be obtained with
very high
accuracy from ratios  of matrix elements of the conserved and
local currents \cite{ma,msv}.
{}From the amputated Green functions
\beq \Gamma_{1}(pa)=\frac{1}{12} {\rm tr} \Bigl( \Lambda_1(pa)\hat 1 \Bigr) ,
\nn \eeq
\beq \Gamma_{\gamma_5}(pa)=\frac{1}{12} {\rm tr} \Bigl( \Lambda_
{\gamma_5}(pa) \gamma_5 \Bigr) , \nn
\eeq
\beq \Gamma_{\gamma_\mu \gamma_5}(pa)=\frac{1}{48} {\rm tr} \Bigl( \Lambda_
{\gamma_\mu \gamma_5}(pa) \gamma_5 \gamma_\mu \Bigr) ,
\eeq
 by imposing the renormalization conditions (\ref{eq:zdef}),
at several
values of $p^2=\mu^2$, to be specified in sec.\ \ref{sec:nonper},
we have
determined $Z_{S,P,A}\Bigl(\mu,g(a)\Bigr)$,
as functions of $\mu^2$.

If we wish to compute
  $Z_{V^L}$  using the procedure defined in sec.\ \ref{sec:casivari}
we clearly cannot take $Z_\psi$ from eq.\ (\ref{eq:zpsil}).
{}From the Ward identity we have
\beq Z_\psi = - i \frac{1}{12}
{\rm tr} \Bigl( \frac{\partial S(p)^{-1}}{\partial \pslash}
\Bigr)\vert_{p^2=\mu^2} \label{eq:zpsiq}
\eeq
where $\partial S(p)^{-1}/\partial \pslash = 1/4 \times
\gamma_\rho \partial S(p)^{-1}/\partial p_\rho$.
To avoid derivatives with respect to a discrete variable,
we have used
\beq Z_\psi^{\prime}\Bigl(\mu , g(a) \Bigr)= \frac {{\rm tr } \Bigl(
-i \sum_{\lambda=1,4} \gamma_\lambda \sin( p_\lambda a ) S^{-1}(pa)
\Bigr)}{4 \sum _{\lambda=1,4} \sin^2 p_\lambda a }\vert_{p^2=\mu^2}.
\label{eq:zpsi1} \eeq
In one loop perturbation theory, $Z_\psi^{\prime}=Z_\psi$,
in the Landau gauge\footnote{In general, they  differ  by a finite
term of order $\alpha_s$.}. By using
\beq \Gamma_{\gamma_\mu }(pa)=\frac{1}{48} {\rm tr} \Bigl( \Lambda_
{\gamma_\mu }(pa) \gamma_\mu \Bigr) , \label{eq:gvl}
\eeq
and imposing the renormalization conditions (\ref{eq:zdef}), we have determined
$Z_{V^L}$ at several $\mu^2$. From the general discussion of the previous
section, one expects that $Z_{V^L}$, defined through the present procedure,
 is independent of $\mu^2$, up to
terms of $O(a)$, as  has been effectively found,
see sec.\ \ref{sec:nonper}.

In actual numerical simulations, the  renormalization procedure
could be obscured by lattice artefacts, i.e.\  terms of $O(\mu a)$, which
are present
 when one imposes
the renormalization conditions at large values of $\mu^2$,  required  to avoid
large
non-perturbative or higher  order effects.
 As shown in ref.\ \cite{noi} (we use here the same notation),  in order to
reduce discretization errors, due to the finite value of the
lattice spacing, one can compute correlation functions using a
nearest-neighbour
improved fermion action \cite{slami} \beq
S^{I}_{F}=
S^{W}_{F} +
a^{4} \sum_{x,\mu \nu} \left[-ig_{0}
\frac{ar}{4}\bar{\psi}(x)
\sigma_{\mu \nu}F_{\mu \nu}(x)\psi(x) \right],
\label{eq:si}
\eeq
where $S^{W}_{F}$ is the usual Wilson action, $r$ is the Wilson parameter
and $F_{\mu\nu}$ is the lattice field strength tensor.
Moreover,
it is also necessary to use improved operators
$O_{\Gamma}^{I}$ defined as
\beq
O_{\Gamma}^I(x) =
\bar{\psi}(x) \overleftarrow{R}_{x}
\Gamma
\overrightarrow{R}_{x} \psi(x),
\label{eq:oi}
\eeq
where
the rotations of the fermion fields
are defined as
\beq
\overrightarrow{R}_{x}=
\Bigl(1-a {r\over{4}} (\overrightarrow{\Dslash}_{x} - m_0)\Bigr)
{}~~~,~~~
\overleftarrow{R}_{y}=
\Bigl(1+a {r\over{4}} (\overleftarrow{\Dslash}_{y} + m_0)\Bigr).
\label{eq:rotaz}
\eeq
By using the improved action and operators, one  eliminates discretization
errors of $O(a)$, and only those of  $O(\alpha_s a)$ or $O(a^2)$, or higher,
remain.
In eq.\ (\ref{eq:oi}), $\overrightarrow R_x$ and
 $\overleftarrow R_y$ can be replaced by
\beq
\overrightarrow{R^\prime}_{x}=
1-a {r\over{2}} \overrightarrow{\Dslash}_{x}
{}~~~,~~~
\overleftarrow{R^\prime}_{y}=
1+a {r\over{2}} \overleftarrow{\Dslash}_{y},
\label{eq:rotaz1}
\eeq
in computations of  the hadron spectrum and of on-shell
operator matrix  elements. As  discussed in refs.\ \cite{tass2,newi},
 the improvement procedure is equivalent to expressing all correlation
functions
in terms of  ``effective'' quark propagators
\beq
S^{eff}_i(x \vert 0)=
\Bigl(1-a{r\over{2}}
\overrightarrow{\Dslash}_{x} \Bigr)
S^{I}_i(x \vert 0)
\Bigl(1+a{r\over{2}}
\overleftarrow{\Dslash}_{y=0} \Bigr)
\label{eq:seff}
\eeq
where $S^{I}_i(x \vert 0)$ is the fermion propagator of the improved theory
from the point $x$ to the origin.
In order to eliminate  terms of $O(a)$ in the
off-shell correlation functions, used in this study, a further step is
necessary
\cite{wi}:
we have to shift the quark propagator by a local term, which in Fourier space
appears as a constant, to be added to its diagonal components
\beq S^{AB}_{\alpha \beta}(pa) = S^{eff,AB}_{\alpha \beta}(pa)
 + \frac{r}{2}
\delta^{AB}
\delta_{\alpha \beta} \nn \eeq

In the following, the vector current  $V^L_\mu$,  the axial vector current
 $A^L_\mu$, the pseudoscalar density $P$
and the scalar density  $S$ are improved operators,
obtained  by using eq.\ (\ref{eq:seff}).
